\documentclass[11pt]{article}

% Shared notes settings for Elements of AIML lectures (UPES).
% Keep this file free of \documentclass to allow reuse via \input.

\usepackage[utf8]{inputenc}
\usepackage[T1]{fontenc}
\usepackage[a4paper,margin=1in]{geometry}
\usepackage{hyperref}
\usepackage{xurl}
\usepackage{booktabs}
\usepackage{amsmath}
\usepackage{amssymb}
\usepackage{listings}
\usepackage{xcolor}
\usepackage{enumitem}
\usepackage{parskip}

% Code listings (general-purpose).
\lstset{
  basicstyle=\ttfamily\small,
  keywordstyle=\color{blue},
  commentstyle=\color{gray},
  stringstyle=\color{teal},
  showstringspaces=false,
  columns=fullflexible,
  frame=single,
  framerule=0.2pt,
  breaklines=true
}

\setlist[itemize]{topsep=0.3em,itemsep=0.2em}
\setlist[enumerate]{topsep=0.3em,itemsep=0.2em}

% Simple per-section source line.
\newcommand{\notesources}[1]{\par\noindent{\small\color{gray}\textit{Sources:} \sloppy #1\par}}

% Unit-wise source strings for this subject (used by slides/notes).
% Path is relative to the lecture `latex/` folder (the compilation CWD).
\IfFileExists{../../../_shared/aiml-sources.tex}{%
  % Source strings for Elements of AIML (used by slides + notes).

\newcommand{\AIMLRepoURL}{https://github.com/tali7c/Elements-of-AIML}

\newcommand{\AIMLLabSources}{%
Provost and Fawcett, \textit{Data Science for Business}.;%
McKinney, \textit{Python for Data Analysis}.;%
WEKA Documentation (Waikato Environment for Knowledge Analysis), accessed 2026-02-28.;%
Matplotlib and Seaborn documentation, accessed 2026-02-28.%
}

\newcommand{\AIMLUnitISources}{%
Rich and Knight, \textit{Artificial Intelligence}, McGraw Hill, 2017.;%
Russell and Norvig, \textit{Artificial Intelligence: A Modern Approach} (4e), Ch. 1--2 (Introduction, Intelligent Agents).;%
Poole and Mackworth, \textit{Artificial Intelligence: Foundations of Computational Agents} (2e), selected sections.%
}

\newcommand{\AIMLUnitIISources}{%
Russell and Norvig, \textit{Artificial Intelligence: A Modern Approach} (4e), Ch. 7--9 (Logical Agents, First-Order Logic, Inference).;%
Huth and Ryan, \textit{Logic in Computer Science} (2e), selected sections on propositional and first-order logic.;%
Genesereth and Nilsson, \textit{Logical Foundations of Artificial Intelligence}, selected chapters.%
}

\newcommand{\AIMLUnitIIISources}{%
Mueller and Massaron, \textit{Machine Learning for Dummies}, For Dummies, 2016.;%
Mitchell, \textit{Machine Learning}, selected chapters on learning basics and evaluation.;%
Scikit-learn documentation, accessed 2026-02-28.%
}

\newcommand{\AIMLUnitIVSources}{%
Mueller and Massaron, \textit{Machine Learning for Dummies}, For Dummies, 2016.;%
James, Witten, Hastie, and Tibshirani, \textit{An Introduction to Statistical Learning} (2e), selected chapters.;%
Scikit-learn documentation, accessed 2026-02-28.%
}

\newcommand{\AIMLUnitVSources}{%
NIST, \textit{AI Risk Management Framework (AI RMF 1.0)}, 2023.;%
Barocas, Hardt, and Narayanan, \textit{Fairness and Machine Learning} (online book), selected sections.;%
Knaflic, \textit{Storytelling with Data}, selected chapters on communicating results.%
}
%
}{%
  % Faculty copies live under `private/` so their compile CWD is different.
  \IfFileExists{../../../../public/_shared/aiml-sources.tex}{%
    % Source strings for Elements of AIML (used by slides + notes).

\newcommand{\AIMLRepoURL}{https://github.com/tali7c/Elements-of-AIML}

\newcommand{\AIMLLabSources}{%
Provost and Fawcett, \textit{Data Science for Business}.;%
McKinney, \textit{Python for Data Analysis}.;%
WEKA Documentation (Waikato Environment for Knowledge Analysis), accessed 2026-02-28.;%
Matplotlib and Seaborn documentation, accessed 2026-02-28.%
}

\newcommand{\AIMLUnitISources}{%
Rich and Knight, \textit{Artificial Intelligence}, McGraw Hill, 2017.;%
Russell and Norvig, \textit{Artificial Intelligence: A Modern Approach} (4e), Ch. 1--2 (Introduction, Intelligent Agents).;%
Poole and Mackworth, \textit{Artificial Intelligence: Foundations of Computational Agents} (2e), selected sections.%
}

\newcommand{\AIMLUnitIISources}{%
Russell and Norvig, \textit{Artificial Intelligence: A Modern Approach} (4e), Ch. 7--9 (Logical Agents, First-Order Logic, Inference).;%
Huth and Ryan, \textit{Logic in Computer Science} (2e), selected sections on propositional and first-order logic.;%
Genesereth and Nilsson, \textit{Logical Foundations of Artificial Intelligence}, selected chapters.%
}

\newcommand{\AIMLUnitIIISources}{%
Mueller and Massaron, \textit{Machine Learning for Dummies}, For Dummies, 2016.;%
Mitchell, \textit{Machine Learning}, selected chapters on learning basics and evaluation.;%
Scikit-learn documentation, accessed 2026-02-28.%
}

\newcommand{\AIMLUnitIVSources}{%
Mueller and Massaron, \textit{Machine Learning for Dummies}, For Dummies, 2016.;%
James, Witten, Hastie, and Tibshirani, \textit{An Introduction to Statistical Learning} (2e), selected chapters.;%
Scikit-learn documentation, accessed 2026-02-28.%
}

\newcommand{\AIMLUnitVSources}{%
NIST, \textit{AI Risk Management Framework (AI RMF 1.0)}, 2023.;%
Barocas, Hardt, and Narayanan, \textit{Fairness and Machine Learning} (online book), selected sections.;%
Knaflic, \textit{Storytelling with Data}, selected chapters on communicating results.%
}
%
  }{%
    % Source strings for Elements of AIML (used by slides + notes).

\newcommand{\AIMLRepoURL}{https://github.com/tali7c/Elements-of-AIML}

\newcommand{\AIMLLabSources}{%
Provost and Fawcett, \textit{Data Science for Business}.;%
McKinney, \textit{Python for Data Analysis}.;%
WEKA Documentation (Waikato Environment for Knowledge Analysis), accessed 2026-02-28.;%
Matplotlib and Seaborn documentation, accessed 2026-02-28.%
}

\newcommand{\AIMLUnitISources}{%
Rich and Knight, \textit{Artificial Intelligence}, McGraw Hill, 2017.;%
Russell and Norvig, \textit{Artificial Intelligence: A Modern Approach} (4e), Ch. 1--2 (Introduction, Intelligent Agents).;%
Poole and Mackworth, \textit{Artificial Intelligence: Foundations of Computational Agents} (2e), selected sections.%
}

\newcommand{\AIMLUnitIISources}{%
Russell and Norvig, \textit{Artificial Intelligence: A Modern Approach} (4e), Ch. 7--9 (Logical Agents, First-Order Logic, Inference).;%
Huth and Ryan, \textit{Logic in Computer Science} (2e), selected sections on propositional and first-order logic.;%
Genesereth and Nilsson, \textit{Logical Foundations of Artificial Intelligence}, selected chapters.%
}

\newcommand{\AIMLUnitIIISources}{%
Mueller and Massaron, \textit{Machine Learning for Dummies}, For Dummies, 2016.;%
Mitchell, \textit{Machine Learning}, selected chapters on learning basics and evaluation.;%
Scikit-learn documentation, accessed 2026-02-28.%
}

\newcommand{\AIMLUnitIVSources}{%
Mueller and Massaron, \textit{Machine Learning for Dummies}, For Dummies, 2016.;%
James, Witten, Hastie, and Tibshirani, \textit{An Introduction to Statistical Learning} (2e), selected chapters.;%
Scikit-learn documentation, accessed 2026-02-28.%
}

\newcommand{\AIMLUnitVSources}{%
NIST, \textit{AI Risk Management Framework (AI RMF 1.0)}, 2023.;%
Barocas, Hardt, and Narayanan, \textit{Fairness and Machine Learning} (online book), selected sections.;%
Knaflic, \textit{Storytelling with Data}, selected chapters on communicating results.%
}
%
  }%
}


\title{Elements of AIML Lab\\Experiment 06 (After-submission Reference): Missing Values and Outliers}
\author{Tofik Ali}
\date{\today}

\begin{document}
\maketitle
\tableofcontents

\section{How to use this reference}
This reference is meant to be shared \textbf{after you submit} Experiment~06.
It explains \textbf{why} missing values and outliers appear, and shows common handling strategies with clear trade-offs.

\section{Missing values: what they mean}
Missing values are not all the same:
\begin{itemize}[leftmargin=*]
  \item \textbf{MCAR} (missing completely at random): missingness is unrelated to any variable.
  \item \textbf{MAR} (missing at random): missingness depends on observed variables (e.g., income missing more often for certain job types).
  \item \textbf{MNAR} (missing not at random): missingness depends on the missing value itself (hardest case).
\end{itemize}
In practice, you rarely prove the type; you justify a reasonable strategy and document it.

\section{Outliers: error vs rare-but-real}
Outliers can be:
\begin{itemize}[leftmargin=*]
  \item \textbf{errors} (wrong unit, extra zero, data entry mistake), or
  \item \textbf{rare-but-real events} (fraud, extreme weather).
\end{itemize}
Therefore, deleting outliers without reasoning can damage the analysis.

\section{Worked example (end-to-end)}
\subsection{Step 1: load dataset and summarize missingness}
\begin{lstlisting}[language=Python]
import pandas as pd

df = pd.read_csv("processed.csv")  # or your dataset

missing_count = df.isna().sum().sort_values(ascending=False)
missing_pct = (df.isna().mean() * 100).sort_values(ascending=False)

print(missing_count.head(10))
print(missing_pct.head(10))
\end{lstlisting}

\subsection{Step 2: choose a missing-value strategy (with justification)}
\paragraph{Option A: dropping.}
Use when missingness is small and dropping does not remove important groups.
\begin{lstlisting}[language=Python]
df_drop = df.dropna()
\end{lstlisting}

\paragraph{Option B: simple imputation (typical baseline).}
\begin{lstlisting}[language=Python]
df_imp = df.copy()

num_cols = df_imp.select_dtypes("number").columns
df_imp[num_cols] = df_imp[num_cols].fillna(df_imp[num_cols].median())

cat_cols = df_imp.select_dtypes(exclude="number").columns
for c in cat_cols:
    df_imp[c] = df_imp[c].fillna(df_imp[c].mode().iloc[0])
\end{lstlisting}

\paragraph{Option C: scikit-learn imputers (pipeline-friendly).}
This is useful when you later combine preprocessing + modeling.
\begin{lstlisting}[language=Python]
from sklearn.impute import SimpleImputer

imp_median = SimpleImputer(strategy="median")
\end{lstlisting}

\subsection{Step 3: detect outliers (IQR rule)}
For one numeric column \texttt{num\_col}:
\begin{lstlisting}[language=Python]
import numpy as np

x = df_imp["num_col"]   # choose a numeric column
q1, q3 = x.quantile(0.25), x.quantile(0.75)
iqr = q3 - q1
lower, upper = q1 - 1.5 * iqr, q3 + 1.5 * iqr

mask = (x < lower) | (x > upper)
print("Outliers by IQR:", mask.sum())
\end{lstlisting}

\subsection{Step 4: handle outliers (two common options)}
\paragraph{Option A: remove outlier rows (only if justified).}
\begin{lstlisting}[language=Python]
df_no_out = df_imp.loc[~mask].copy()
\end{lstlisting}

\paragraph{Option B: cap outliers (winsorization).}
\begin{lstlisting}[language=Python]
df_cap = df_imp.copy()
df_cap["num_col"] = df_cap["num_col"].clip(lower, upper)
\end{lstlisting}

\subsection{Step 5: compare before/after impact}
\begin{lstlisting}[language=Python]
print("Before:\n", df.describe())
print("After imputation:\n", df_imp.describe())
print("After capping:\n", df_cap.describe())
\end{lstlisting}
Add one plot if possible (boxplot before vs after) to show the effect clearly.

\section{What strong observations look like (examples)}
\begin{itemize}[leftmargin=*]
  \item ``\texttt{Age} had 4.2\% missing values. We used median imputation because the distribution is skewed.''
  \item ``IQR detected 12 extreme values in \texttt{Income}. We capped them to reduce distortion while keeping rare cases.''
  \item ``After capping, max income reduced from 9.8M to 1.2M and mean decreased, indicating extreme values were pulling the mean upward.''
\end{itemize}

\section{Troubleshooting}
\begin{itemize}[leftmargin=*]
  \item \textbf{All values flagged as outliers:} check if the column is categorical coded as numbers, or if units are inconsistent.
  \item \textbf{Imputation changes categories incorrectly:} avoid mode fill when the column has many unique IDs.
  \item \textbf{Leakage warning (for ML):} always fit imputers/outlier rules on training data only when doing model evaluation.
\end{itemize}

\notesources{\AIMLLabSources}
\end{document}

